\documentclass[10pt,a4paper]{amsart}
\usepackage[margin=19mm]{geometry}
\usepackage{amsmath,amssymb,mathtools}
\usepackage[scr=rsfs,cal=euler]{mathalfa}
\usepackage{xfrac}
\usepackage[unicode,hidelinks,bookmarks=false]{hyperref}
\newcommand{\ie}{i.e.\ }
\newcommand{\cf}{cf.\ }
\newcommand{\resp}{respectively\ }
\newcommand{\Subsection}{\subsection}
\providecommand{\nobreakdash}{\nobreak\hbox{-}}
\setlength{\emergencystretch}{2em}
\multlinegap0pt
\usepackage{amssymb}
\newtheorem{lemma}{Lemma}
\newtheorem{proposition}{Proposition}
\def\N{\mathbb{N}}
\def\Om{\Omega} 
\def\R{\mathbb{R}}
\DeclareMathOperator{\Tr}{Tr}
\DeclareMathOperator{\dist}{dist}
\def\k{\kappa} 
\def\subjet{J^{2,-}}
\title{Power-logconcavity of the Laplacian ground state}
\author{Graziano Crasta, Ilaria Fragal\`a}
\date{Source: arXiv:2602.24093v1 (CC BY 4.0)}
\begin{document}
\maketitle

\noindent\emph{Source and licence.} Power-logconcavity of the Laplacian ground state. Version arXiv:2602.24093v1 (CC BY 4.0), distributed by arXiv under the Creative Commons Attribution 4.0 International licence, http://creativecommons.org/licenses/by/4.0/, as recorded in the arXiv licence metadata for this exact version and shown on the abstract page https://arxiv.org/abs/2602.24093v1. Full licence notice: \texttt{licenses/arxiv-2602.24093-CC-BY-4.0.txt}.\par\smallskip

\noindent\textit{Editorial scope.} Counted unit: the article's proposition p:viscosuper (source0.tex lines 588--590). The article's complete original proof of it is delivered with it (source0.tex lines 592--719, one \texttt{proof} environment of 128 lines). No mathematics is written, repaired or re-derived: the statement and the proof are the article's own lines, in the article's own notation, and no retained line is substituted or re-flowed.  Retained uncounted as the source-local context the proof needs: lines 262--269; lines 285--303; lines 422--452; lines 512--514; lines 518--520; lines 551--557.  Nothing was omitted from any retained span, so the package carries no omission record.  No retained line contains a citation, so the wrapper carries no bibliography and no bibliography entry.  No retained line contains a figure environment, an image-inclusion command or drawing code, so no image is delivered and the package declares no asset dependency.  Editorial formatting only, in addition: an AMS \texttt{amsart} preamble that restates the article's own declarations and macros used by the retained lines, namely the environments \texttt{lemma}, \texttt{proposition} on separate counters, together with the standard packages those lines need.  The retained environments print in this document as Lemma 1, Proposition 1 in order of appearance, whereas the article prints the same items with the numbering of its own full document.  The complete native source of the article as distributed in the arXiv e-print for this exact version is bundled unchanged as \texttt{sources/195446/source0.tex}.
 \begin{equation}\label{f:w**} 
w^{**} ( x):= 
\inf \Big \{ \sum_{ i = 1}^{m} t _i w(x_i) \ :\ 
m \in \N,\
t _ i\geq 0,\ x_i \in \Omega\,, \
\sum _{i = 1} ^ m t _ i = 1,\ 
\sum _{i = 1} ^ m t _ i x_i = x  \Big \}.
\end{equation}

\begin{lemma}
\label{l:bdr}
Let $\Omega$ be an open bounded convex subset of $\R^n$,  and 
let $w\colon\Omega\to\R$ be a lower semicontinuous function
satisfying
$w(x) \to +\infty$ as $\dist(x, \partial\Omega) \to 0$.

Then $w^{**}(x) \to +\infty$ 
as $\dist(x, \partial\Omega) \to 0$.

Moreover, for every $x\in\Omega$, there exist
$t_1, \ldots, t_m > 0$, $m\leq n+1$, with $\sum_{i=1}^m t_i = 1$,
and points $x_1, \ldots, x_m \in \Omega$, such that 
\begin{equation}
\label{f:min}
w^{**}(x) = \sum_{i=1}^m t_i w(x_i),
\end{equation}
i.e., the infimum which defines $w ^ {**}$ according to \eqref{f:w**} is attained.
\end{lemma}

\begin{lemma}
\label{l:ALLp1}
Let $\Omega$ be an open, bounded, convex subset of $\R^n$,  and 
let $w\in C^2(\Omega)$ be such that $w(x) \to +\infty$ as $\dist(x, \partial\Omega) \to 0$.

Given $x\in\Omega$, let $t_1, \ldots, t_m > 0$, $m\leq n+1$, with $\sum_{i=1}^m t_i = 1$,
and let $x_1, \ldots, x_m \in \Omega$ be optimal for the infimum defining $w ^ {**} (x)$ according to \eqref{f:w**}.

If $(p,A) \in \subjet   w^{**}(x)$,  then 
the following properties hold:
\begin{itemize}
\item[(i)] $\nabla w(x_i) = p$
for every $i = 1,\ldots,m$;
\smallskip

\item[(ii)] the matrices $A_i := \nabla ^2 w(x_i)$ are positive semidefinite for every $i = 1,\ldots,m$;
\smallskip

\item[(iii)] if all the matrices $A_i$, $i=1,\ldots,m$, are positive definite,
then
\[
A \leq \left(\sum_{i=1}^m t_i A_i^{-1}\right)^{-1},
\qquad
\Tr(A) \leq 
\left(\sum_{i=1}^m t_i \Tr (A_i^{-1})\right)^{-1};
\]

\item[(iv)] if at least one of the matrices $A_i$ is degenerate, then
$\Tr(A) \leq 0$.
\end{itemize}
\end{lemma}

\begin{equation}\label{f:norm} 
w  _\k:= (- \log (\k u) ) ^ { 1/2} \qquad \text{ in } \Om\,. 
\end{equation}   

\begin{equation}\label{f:bvw}
 - \Delta w + \frac{1}{w} \big [ ( 2 w ^ 2 -1 ) |\nabla w| ^ 2 +\frac{\lambda _ 1 (\Om) }{2} \big ] = 0 \qquad \text{ in } \Om\,. 
 \end{equation} 

\begin{proposition}\label{p:key}  Let $\k \in (0, 1)$ and assume that, at some point $x \in \Om$, one has $w _\k   ^ { **}  (x) < w _\k  (x) $. 
Let $x _i \in \Om$, $i = 1, \dots, m$  (with $m \geq 2$) be points where the infimum which defines   $w _\k   ^ { **}  (x)$ according to
\eqref{f:w**} is attained, and let $p_\k$ 
denote  the common value of  $\nabla w _\k$ at any of the points   $x_i$, as given by Lemma  \ref{l:ALLp1}  (i). 
Then 
 $$|p_\k | ^ 2 \geq  \frac{\pi ^ 2}{2 D _\Om ^ 2} \,.$$   
\end{proposition} 

\begin{proposition}\label{p:viscosuper}   
For any $\k \in (0, \overline k (\Om)]$, the convex envelope $w  _\k ^{**}$ is a viscosity  supersolution  to   equation \eqref{f:bvw}. 
  \end{proposition} 

\begin{proof}
Let $\k \in (0, \overline k (\Om)]$ be fixed. We have to show that, setting  
$$F (s , \xi, X):= - {\rm {\rm Tr} } ( X)  + \frac{1}{s} \big [ ( 2 s ^ 2 -1 ) |\xi| ^ 2 +\frac{\lambda _ 1 (\Om) }{2} \big ] \,,$$ 
it holds that
\begin{equation}\label{f:goal1}
F ( w _\k ^{**}(x), p, A ) \geq 0 \qquad \forall x \in \Om \,, \quad \forall (p, A) \in \subjet w _\k  ^ { ** } ( x) \,.
\end{equation} 


Let $x \in \Om$ be fixed. 
We first observe that \eqref{f:goal1} holds trivially  if $w _\k ^{**}(x) = w _\k ( x)$. Indeed in this case we have  
$ \subjet w _\k  ^ { ** } ( x) \subseteq \subjet w _\k   ( x)$, so that 
$(p, A) \in \subjet w _\k   ( x)$, and thus \eqref{f:goal1} holds because $w _\k$ is a classical solution, hence a viscosity solution,  to equation \eqref{f:bvw}. 
Observe that this case happens, in particular, when $p=0$, because in this case $x$ must coincide with the unique minimum point of $w_\k$,
i.e., the unique maximum point of $u$.

Thus, in the remaining of the proof we are going to assume that $w _\k ^{**}(x) < w _\k ( x)$ and $p\neq 0$. 



By Lemma \ref{l:bdr}, there exist $t ^\k  _i \in (0, 1)$ and  $x ^\k _i \in \Om$, $i = 1, \dots, m_\k$ (with $m _\k \geq 2$), 
such that
\begin{equation}\label{f:env}
\sum _{i=1}^{m _\k} t ^\k  _i =1\,, \qquad
x = \sum _{i=1}
 ^{m _\k} t ^\k  _i x ^\k  _i\,, \qquad w^{**} _\k  (x)  = \sum _ { i = 1} ^{ m_\k} t ^\k   _i w _\k  ( x ^\k _i) \,.
\end{equation} 
 
Throughout the proof, the parameter
$\k \in (0, \overline \k (\Om)]$, the point $x \in \Om$, and accordingly the
numbers $t^\k _i $ and  the points $x^\k _i $ in \eqref{f:env}  will remain fixed. 
Taking also into account that,  by Lemma \ref{l:ALLp1} (ii), the vector  $\nabla  w  _\k ( x_i)$  is  independent of  $i \in \{ 1, \dots, m _\k\}$,
we are going to write for brevity  
$$w ^{**} := w^{**} _\k  (x)$$
and, for every $i = 1, \dots, m: = m _\k$ (with $m \geq 2$), 
$$t_i:= t ^ \k _ i \, , \qquad w^i:= w  _\k ( x_i),    \qquad p:= \nabla  w  _\k ( x_i)\ , \qquad A _{ i}: = \nabla ^ 2 w _\k  ( x_i) \,.$$ 
Moreover, we set 
 \begin{equation}\label{f:defa} 
 a:=  \frac{\lambda _1 (\Om)}{4| p | ^ 2 } - \frac 1 2 \,,
 \end{equation}
so that
$$\begin{aligned}
& F (w^{**}, p, A )  = - \Tr (A)+  2 |p| ^ 2 \frac{(w^{**} ) ^ 2 + a}{w ^ {**}}\,,
\\ 
& F (w^i, p, A_i )  =- \Tr (A_i) + 2 |p| ^ 2 \frac{(w^{i} ) ^ 2 + a}{w ^ {i}} \,.  
 \end{aligned}$$
  
 
By Lemma \ref{l:ALLp1} (ii), we know that all the matrices $A_i$ are positive semi-definite.  
  Let us first prove the inequality \eqref{f:goal1} when  all of them are 
positive definite. 
In this case,  by Lemma \ref{l:ALLp1}  (iii),  it holds that
 $$\Tr(A) \leq 
\left(\sum_{i=1}^m t_i \Tr (A_i^{-1})\right)^{-1}\,. $$ 
 Hence,  $$\begin{aligned} 
& F (w^{**}, p, A )  
 =  - {\rm Tr} ( A)  +  2 |p| ^ 2 \frac{(w^{**} ) ^ 2 + a}{w ^ {**}}   \\ 
 & \geq   - \Big \{  \sum _ { i = 1} ^ m  t _i  \big [  {\rm Tr} (A _i )   \big ]   ^ { -1}   \Big \} ^ { -1}  
 + 2 |p| ^ 2 \frac{(w^{**} ) ^ 2 + a}{w ^ {**}}    \\ 
 & =   - \Big \{  \sum _ { i = 1} ^ m  t _i  
 \Big [  2 |p| ^ 2 \frac{(w^{i} ) ^ 2 + a}{w ^ {i}}   \Big ]      ^ { -1}   \Big \} ^ { -1}   + 2 |p| ^ 2 \frac{(w^{**} ) ^ 2 + a}{w ^ {**}}     \,,
   \end{aligned} 
 $$ 
 where in the last equality we have used the identity
  \begin{equation}\label{f:equazz} 
   F (w^i, p, A_i ) = 0 \qquad \forall i = 1, \dots, m, 
 \end{equation}
  holding since $w _\k $ satisfies the pde  \eqref{f:bvw}
 at the points $x _i$'s.   
 
 
 We infer that the target inequality \eqref{f:goal1} is fulfilled provided
 \begin{equation}\label{f:goal2} \frac{w^{**}}{(w^{**}) ^ 2  + a } 
 \leq \sum _ { i = 1} ^ m  t _i  
 \frac{w^i} { (w^{i} ) ^ 2  + a   }  \,.
 \end{equation} 
We now observe that, since we are assuming that the matrices $A_i$ are positive definite, and since
the equalities \eqref{f:equazz}  hold, we have 

 $$0 < {\rm Tr} ( A_i) = 2 |p | ^ 2 \, \frac {(w^ {i} ) ^ 2  + a  }{w^ i}   \qquad \forall i = 1, \dots, m.$$ 
 We deduce that 
 \begin{equation}\label{f:app} 
 w ^ i \in \mathcal R:=  \big \{  s \in (0,   + \infty)  \text{ such that } s ^ 2 > - a  \big \} \qquad \forall i = 1, \dots, m \,.
 \end{equation} Hence, we see that
the inequality \eqref{f:goal2} is satisfied provided  the values $\{w^ 1, \dots, w ^m, w ^{**}\}$ fall in the subregion $\mathcal R_0$ of $\mathcal R$ where the function 
 $$\varphi _ a ( s):= \frac{s}{ s ^ 2 + a  }\,,$$ 
 
is convex. 
 Since 
 $$\varphi _ a''  ( s):= \frac{2s (s^2- 3a)}{ (s ^ 2 + a ) ^ 3  } \,,$$ 
 we have
 $${\mathcal R}_0=  \mathcal R \cap \{ s^2 \geq 3 a \}  \,.$$ 
 So we distinguish two cases:
 
 \smallskip
 {\it Case 1}:  $a \leq 0$. In this case we have 
 $${\mathcal R}_0=  \mathcal R = \{ s > \sqrt { - a \}}\,. $$ 
 By \eqref{f:app}, we have that, for every $i = 1, \dots, m$,  the values $w ^ i$ lie in $\mathcal R _0$; by the convexity of $\mathcal R _0$, also  $w ^{**}$ lie in $\mathcal R _0$. Thus
 the required inequality \eqref{f:goal2} is satisfied, and our proof is achieved.
 
 \smallskip 
  {\it Case 2}:  $a >0$. In this case,  imposing that $w ^ i \in \mathcal R _0$ amounts to ask that 
\begin{equation}\label{f:stimap}
|p| ^2\geq \frac{\lambda _ 1 (\Om) }{4} \frac{1}{\frac{ (w ^ i) ^ 2}{3} + \frac {1}{2}  } \qquad \forall i = 1, \dots, m \,.
\end{equation}
We claim that these inequalities are fulfilled thanks to the assumption that the constant $\kappa$ in \eqref{f:norm} does not exceed $\overline \kappa (\Om)$. 
Indeed, since $m \geq 2$, we are in a position to apply Proposition \ref{p:key}: it implies that
a sufficient condition for the validity of the inequalities in  \eqref{f:stimap} is 
 $$ \frac{\pi ^ 2}{2 D _\Om ^ 2} \geq   \frac{\lambda _ 1 (\Om) }{4} \frac{1}{\frac{ (w ^ i ) ^ 2}{3} + \frac {1}{2}  } \qquad \forall i = 1, \dots, m \,, $$ 
 or equivalently 
 $$( w^i ) ^ 2 =  - \log (\k u(x_i)) )   \geq  \frac 3 2 \Big ( \frac{ \lambda _ 1 (\Om) D _\Om^2  }{\pi ^ 2}  -1 \Big )  \qquad \forall i = 1, \dots, m   \,.$$ 
Such inequalities are satisfied  because  $0 < u (x_i) \leq 1$ and  $\kappa  \leq \overline \kappa (\Omega)$.  

 \smallskip
It remains to prove that  the inequality 
\eqref{f:goal1}  is satisfied also when some of the $A_i$ is degenerate. In this case, by  Lemma \ref{l:ALLp1} (iv), we know that ${\rm Tr} (A) \leq 0$. Hence,  
 $$
F (w^{**}, p, A ) 
 =  - {\rm Tr} ( A)  +  2|p|^ 2 \frac{( w^ {**} ) ^2 +a   }{w^{**}  }  \geq  2|p|^ 2 \frac{( w^ {**} ) ^2 +a   }{w^{**}  } \,.
$$
To conclude we observe that: if $a \geq 0$, we have immediately   $
 F (w^{**}, p, A ) \geq 0$; on the other hand, if $a <0$, we have that 
 the map $s \mapsto \frac{ s ^ 2 + a}{s} = s - \frac{|a|}{s}$ is concave in $(0, + \infty)$, and hence
 \[
 F (w^{**}, p, A )\geq  2|p|^ 2 \frac{( w^ {**} ) ^2 +a   }{w^{**}  }  \geq 2|p|^ 2  \sum _ {i= 1} ^m t_i \frac{ (w^ i )  ^ 2 + a}{w^ i}  =    \sum _ {i= 1} ^ m  t_i \Tr (A_i)  \geq 0 \,.
 \qedhere
\]
\end{proof}

\end{document}
